\begin{viuabstract}{Resumen}
Este TFM estudia el transporte cooperativo de cargas heterogéneas mediante tres
subproblemas: SP1 forma coaliciones y asigna roles; SP2 aproxima, acopla,
transporta y sustituye miembros; SP3 coordina rutas y reservas entre misiones.
El modo primario, \emph{Cargo}, trata la carga soportada como un cuerpo rígido
planar.

La propuesta combina referencias LP/MILP, juegos potenciales, dinámica
primal--dual, cierre entero y contratos de recursos, contacto, movimiento,
recuperación y red. Se prueban regulación de cuotas, estabilidad local de Cargo
bajo contacto fijo y realización exacta del \emph{wrench}, terminación de
mejoras de rutas e imposibilidad de garantizar optimalidad global cuando la
desconexión oculta instancias localmente indistinguibles con óptimos incompatibles. Las campañas usan mundos pareados, semillas explícitas y oráculos
globales identificados como techos experimentales.

SP1 regula cuotas bajo condiciones explícitas, separa la capacidad heterogénea
y reduce falsos factibles nominales en la interfaz SP1--SP2. SP2 revela un intercambio
adverso entre seguridad y progreso y restaura los certificados recuperables del
banco estático; en su juego de reparación, el peor equilibrio no admite una
cota uniforme. En SP3, la reserva reduce
interbloqueos en el catálogo discreto; el piloto congestionado no entrega.
Cargo completa casi todas las misiones sin ejecutar como una sola cadena todos
los mecanismos teóricos.

La contribución es una arquitectura verificable de contratos y juegos locales,
más una composición funcional planar. No demuestra estabilidad híbrida global,
optimalidad social, coordinación completamente distribuida ni validez
industrial. El contacto instrumentado, el hardware y \emph{caging} quedan fuera
de la evidencia confirmatoria.

\keywords{Palabras clave}{coordinación multi-robot,
  robots móviles autónomos, formación de coaliciones,
  juegos potenciales, transporte cooperativo}
\end{viuabstract}

\clearpage
